\documentclass[12pt]{article}

\usepackage[dvips,letterpaper,margin=0.75in,bottom=0.5in]{geometry}
\usepackage{cite}
\usepackage{slashed}
\usepackage{graphicx}
\usepackage{amsmath}
\usepackage{amssymb}
\usepackage{braket}
\usepackage[american,fulldiode]{circuitikz}

\begin{document}
\newcommand{\ihbar}{\ensuremath{i \hbar}}
\newcommand{\dPsidt}{\ensuremath{ \frac{\partial \Psi}{\partial t} }}
\newcommand{\dPsidx}{\ensuremath{ \frac{\partial \Psi}{\partial x} }}
\newcommand{\ddPsidx}{\ensuremath{ \frac{\partial^2 \Psi}{\partial x^2} }}
\newcommand{\dPssdt}{\ensuremath{ \frac{\partial \Psi^*}{\partial t} }}
\newcommand{\dPssdx}{\ensuremath{ \frac{\partial \Psi^*}{\partial x} }}
\newcommand{\ddPssdx}{\ensuremath{ \frac{\partial^2 \Psi^*}{\partial x^2} }}

\newcommand{\dphidt}{\ensuremath{ \frac{d \phi}{dt} }}
\newcommand{\dpsidx}{\ensuremath{ \frac{d \psi}{dx} }}
\newcommand{\ddpsidx}{\ensuremath{ \frac{d^2 \psi}{dx^2} }}


\date{\vspace{-5ex}}

\title{Homework Assignment 5 \\ Toy Model }

\maketitle

\noindent
All problems are graded on effort only.\\

\noindent
{\bf Griffiths: 6.8, 6.12, 6.13, 6.15}\\[5pt]


\noindent
{\bf Note on P6.8:} There is a typo in the answer in some versions of the {\bf digital} textbook for the course.  The correct answer is:
$$\frac{\partial \sigma}{\partial \Omega} = \left( \frac{\hbar}{8 \pi m_b c} \right)^2 | \mathcal{M} |^2$$


\noindent
{\bf Hints on P6.15:} This is a challenging problem, so I will give you plenty of help.\\[5pt]

\noindent
Let's use a consistent labeling to keep things clear below.  Incoming particles: $A$ has four-momentum $p_a$, $B$ has $p_b$.  Outgoing particles: $A$ has $p_1$ and $B$ has $p_2$.  The scattering angle $\theta$ is the angle between the incoming and outgoing $A$, so:
$$\vec{p_a} \cdot \vec{p_1} = |\vec{p_a}| |\vec{p_1}| \cos\theta$$\\

\noindent
{\bf P6.15a}: The outgoing particles are not identical, so the second
diagram does not come from crossing the outgoing particles (as in $A+A
\to B+B$).  This process still has two diagrams, one with the internal
$C$ line horizontal (we call this a $s$-channel process) and one with
internal $C$ line vertical (we call this a $u$-channel process).  Your
answer will look like this:
$${\cal M} = g^2 \left[\frac{1}{(X-m_C^2c^2} + \frac{1}{(Y-m_C^2c^2} \right]$$
where $X$ and $Y$ are each one of the Mandelstrom variables $s$,$t$,or $u$:
\begin{eqnarray*}
s &=& (p_a + p_b)^2 \\
t &=& (p_a - p_1)^2 \\
u &=& (p_a - p_2)^2 \\
\end{eqnarray*}
~\\[5pt]

\noindent
{\bf P6.15b}: Notice that now $\vec{p_a}=-\vec{p_b}$ and $\vec{p_2}=-\vec{p_1}$.  Also:
$$|\vec{p_a}| = |\vec{p_b}| = |\vec{p_1}| = |\vec{p_2}| \equiv p$$  First show that:
$$s = (p_a+p_b)^2 = 4 E^2$$
and
$$u = (p_a-p_2)^2 = -2 p^2 \, (1 + \cos\theta) = -4 p^2 \cos^2\left( \frac{\theta}{2}\right)$$
Make sure your answer for $d\sigma/d\Omega$ is a function of $E$ and $\theta$ variables only.  If you like, you can simplify a factor like:
$$\left( \frac{a}{b} + \frac{c}{d}\right)^2 \to \left(\frac{ad+cd}{bd}\right)^2$$
but it won't get that much prettier.\\[5pt]

\noindent
{\bf P6.15c:} Use P6.8.  Also show that in the limit of large $m_B$:
$$s = u = m_B^2 c^4$$

\noindent
{\bf P6.15d:} Easy integral!  No $\theta$ or $\phi$ dependence!

\end{document}




